\begin{abstract}
NISAR's dual-band L-band interferograms come with a split-spectrum ionospheric phase screen, which studies of large
earthquakes subtract. Near the surface ruptures of the 2026 Venezuela and Panama earthquakes, the screens in NISAR's
PR GUNW products contain a scaled, sign-flipped copy of the coseismic deformation. Within 10\,km of the trace their
15\,km high-pass regresses on the data's at $-0.02$ to $-0.12$ ($r$ down to $-0.82$).

We locate where the deformation enters by running the two-band, the reformulated and a reordered split-spectrum
estimate, with and without ISCE3's whole-cycle corrections. All run on identical cells, with an identical filter and
an identical main-band unwrapping, on synthetic ruptures with known truth and on five NISAR tracks. Three mechanisms
appear.
\begin{itemize}
  \item The correction passes read the deformation's curvature as unwrapping errors. On synthetics they copy
        4.5--4.8\% of the high-passed deformation at every slip, 6\% without a fault-zone exclusion.
  \item Unwrapping a field that carries deformation fails once the rupture splits the scene and the field's share of
        the step exceeds about $\pi$. That field is the side band, or the double difference with its 4.4\% share. On the
        Panama pair it puts the cut-off block one cycle of the double difference, 1.4\,m of LOS, off.
  \item The product's own main-band unwrapping, 25--32 cycles off across the Panama rupture, accounts for most of
        that product's copy.
\end{itemize}
On the Venezuela tracks the corrections, run on matched cells, reproduce 25--52\% of the products' leak beyond our
own; the rest is not attributed here.

The reordered estimate removes the scaled, unwrapped main band from the wrapped side band before anything is
unwrapped. It is algebraically identical to the reformulated estimate of Wegm\"uller et al.\ (2018) and to the one
ISCE3 computes. But the only field it unwraps holds no deformation, so it needs no cycle corrections. On every
synthetic, its error 2--5\,km from the trace stays at its zero-slip level of 8--12\,mm of LOS, which is the filter's.
The GUNW recipe's error is 47\,mm at 4\,m of slip, and the forms that unwrap deformation fail by metres where the
rupture splits the scene. The reordered estimate still needs the main band unwrapped correctly, at 3.21\,rad of
screen per cycle.

That same 3.21\,rad makes the side band a tie for a block that a cut along the rupture isolates: the ionosphere is
continuous across the gap, so the raw estimate's step there reads the block's integer. On fake boundaries in eight
NISAR scenes, no tie, the median over a boundary's 20\,km windows, reads a wrong integer; across a Venezuela rupture
whose integer is known, it reads 0.02. On three of the four Panama pairs it finds the pipelines' range-offset ties
one cycle short, which had left a 5--6\,cm fake step in their screens within 8\,km of the trace. The offsets themselves,
read the same way, agree with it; comparing two far-field planes at the trace was what had moved the ties.

On the synthetics, three parts of the estimator set its zero-slip error near the trace. The cells sum steep fringes
differently in the two bands, the excluded zone is filled with a plane wider than the zone, and the filter's
inverse-variance weighted mean lags the ionosphere's gradient across the damage zone. Taking the main band out of every
pixel, a fill matched to the gap and a weighted local plane as the final step bring that error from 16 to 2\,mm within
2\,km of the trace. On real data they are millimetre corrections. On 34 fake gaps in eight NISAR screens the fills
differ by about a millimetre. Next to the D141 rupture the screen is limited instead by the main band's local integers:
two unwrappings disagree on 18--78\% of the cells within 6\,km north of the trace, and the narrower recipes follow
those cells where the production screen's wide fill averages them away.

Over ice the same failures appear without an earthquake. On two Antarctic pairs whose PR screen Joughin (2026) found
offset across a glacier's shear margin, that screen is ours plus whole cycles:
\begin{itemize}
  \item the main band's split between its components, 9 and 8 cycles at 3.21\,rad each;
  \item one cycle of the double difference over a 12\,000\,km$^2$ block;
  \item about two main-band cycles that the correction passes moved over a fast tributary.
\end{itemize}
Where the screens differ beyond those integers, the range offsets side with ours. Across a 160\,rad step the side band
and the range offsets tie the components alike, to 0.05 cycles, where velocity maps read one or two cycles low.
\end{abstract}
